\documentclass[11pt]{article}
\usepackage[margin=1in]{geometry}
\usepackage{times}
\usepackage{latexsym}
\usepackage{amsmath}
\usepackage{amssymb}
\usepackage{booktabs}
\usepackage{multirow}
\usepackage{graphicx}
\usepackage{xcolor}
\usepackage{array}
\usepackage[T1]{fontenc}
\usepackage[utf8]{inputenc}

\newcommand{\balacc}{\textsc{BalAcc}}
\newcommand{\fpr}{\textsc{FPR}}
\newcommand{\prec}{\textsc{Prec}}

\title{Layer Selection Bias Reverses Conclusions in Hidden-State Safety Monitoring:\\A Systematic Evaluation Across Four Reasoning Models}

\author{Anonymous}

\begin{document}
\maketitle

\begin{abstract}
Reasoning language models can self-jailbreak during extended chain-of-thought generation, progressively abandoning safety reasoning until a \emph{commitment point} after which harmful content is produced.
Detecting this commitment in real time is an open problem, yet prior evaluations of hidden-state (HS) monitoring have not examined how the choice of layer selection metric affects conclusions.
We evaluate seven detection methods---temporal logistic regression, static probes, GRU sequence models, Mahalanobis distance, BOCPD/CUSUM change-point detection, and text-based probes---across four reasoning models spanning two architectures (Llama, Qwen), two training paradigms (R1-distilled, STILL), and two scales (7--8B, 32B).
We discover that the composite metric \texttt{bal\_acc $\times$ crossing\_rate}, used for layer selection, systematically biases toward shallow layers in 32B models, entirely masking HS probe advantages and producing false-negative conclusions about hidden-state monitoring efficacy.
After correcting for this bias via a multi-metric audit protocol, HS probes significantly outperform text probes on all four models (12/12 Holm--Bonferroni tests, precision $+$15--24pp).
However, no method achieves deployment-ready performance: safe-trace false positive rates range from 11\% to 52\%, cross-model transfer fails completely (best 67.3\% vs.\ 50\% chance), and 0/4 models simultaneously meet practical thresholds for precision ($>$70\%), coverage ($>$60\%), and false positive rate ($<$20\%).
We release our multi-metric sensitivity analysis framework and argue that future hidden-state safety evaluations must report results under multiple layer selection criteria to avoid systematic bias.
\end{abstract}


\section{Introduction}

Recent reasoning-capable language models---such as DeepSeek-R1 \cite{deepseekai2025deepseekr1} and QwQ \cite{qwq32b}---generate extended chains of thought (CoT) that can span hundreds of tokens. A growing body of work documents \emph{self-jailbreaking}, where the model's own reasoning process progressively erodes safety guardrails, culminating in a \emph{commitment point} after which harmful content is produced \cite{wen2025thinkingmakesllmssusceptible,zhou2025hiddenriskslargereasoning}.

Monitoring these internal deliberations for safety commitment is critical for deployment. Two monitoring modalities are available: (1) \textbf{text-based} monitoring of the generated CoT tokens, requiring no model internals; and (2) \textbf{hidden-state (HS)} monitoring of intermediate representations, requiring white-box access. Prior work on hidden-state safety detection has focused on adversarial input attacks against instruction-tuned models \cite{wang2024mitigating,ji2024llmguard}, leaving the reasoning-model self-jailbreak setting unexplored.

We provide the first systematic comparison of detection methods for reasoning-model safety commitment. Our evaluation spans four models (DeepSeek-R1-Distill-Qwen-32B, QwQ-32B, DeepSeek-R1-Distill-Llama-8B, OpenThinker-7B), seven detection methods, and five layer selection metrics. Our central discovery is methodological:

\begin{itemize}
\item The composite metric \texttt{bal\_acc $\times$ crossing\_rate} systematically selects shallow layers for 32B models (3--16\% depth vs.\ 94--98\% for corrected metrics), because shallow layers exhibit higher crossing rates due to compressed representations. This bias entirely masks the HS advantage at 32B scale.
\item Under corrected layer selection, HS probes significantly outperform text probes on \emph{all four} models (12/12 Holm--Bonferroni, precision $+$15--24pp).
\item Despite this statistical advantage, \emph{no evaluated method} reaches deployment-ready performance. Safe-trace false positive rates are 39\% at the median, coverage at high-precision operating points drops below 50\%, and probe transfer across models fails completely.
\end{itemize}

These findings carry two implications for the field. First, hidden-state safety evaluations that select a single layer via a single metric risk systematic bias; we propose a multi-metric audit as a minimal safeguard. Second, the HS precision advantage of $+$15--24pp over text must be weighed against the engineering cost of white-box access and per-model probe training---a cost-benefit analysis we provide in \S\ref{sec:discussion}.


\section{Related Work}

\paragraph{Hidden-state safety monitoring.}
TrajGuard \cite{wang2024mitigating} uses Mahalanobis distance on hidden states to detect adversarial attacks with $>$95\% accuracy on standard instruction-tuned models. SafeSwitch and SIREN \cite{ji2024llmguard} extend this to token-level intervention. These methods target adversarial \emph{input} attacks; we evaluate whether they generalize to the reasoning-model self-jailbreak setting where the threat originates in the model's own CoT.

\paragraph{Reasoning model safety.}
Self-Jailbreaking \cite{wen2025thinkingmakesllmssusceptible} and HarmThoughts \cite{zhou2025hiddenriskslargereasoning} document the phenomenon of progressive safety erosion in CoT reasoning. SafeThink and SafeRemind \cite{zheng2025safethink} propose text-level interventions but remain vulnerable to pathological CoT. None provide systematic detection evaluations comparing HS and text modalities.

\paragraph{Representation probing.}
Linear probing of internal representations for behavioral properties is well-established \cite{li2024inferencetime,burns2023discovering}. RepE \cite{zou2023representation} uses mean-difference directions. Recent work shows that probe performance is sensitive to layer choice \cite{belrose2023eliciting}, but no prior study has documented systematic \emph{metric-dependent} layer selection bias.


\section{Method}

\subsection{Task: Safety Commitment Detection}

Given a reasoning trace $\mathbf{x} = (x_1, \ldots, x_T)$ containing an annotated commitment point $\tau$, the detection task is to classify each step $t$ as pre-commitment ($t < \tau$) or post-commitment ($t \geq \tau$). A detector produces a binary prediction per step; trace-level detection requires $K \geq 5$ consecutive positive predictions.

\subsection{Models}

We evaluate four models spanning a $2 \times 2$ design of architecture and scale:

\begin{table}[h]
\centering
\small
\begin{tabular}{llllr}
\toprule
\textbf{Model} & \textbf{Arch} & \textbf{Training} & \textbf{Scale} & \textbf{$n$} \\
\midrule
R1-8B & Llama & R1-distilled & 8B & 43 \\
OT-7B & Qwen & STILL & 7B & 47 \\
QwQ-32B & Qwen & native & 32B & 19--20 \\
R1-32B & Qwen & R1-distilled & 32B & 20 \\
\bottomrule
\end{tabular}
\caption{Model characteristics and test set sizes.}
\label{tab:models}
\end{table}

\subsection{Detection Methods}

We evaluate seven detection approaches:

\begin{enumerate}
\item \textbf{HS Static LR}: Logistic regression on hidden states at a single time step.
\item \textbf{HS Temporal LR}: Logistic regression on concatenated hidden states across a window of $W$ steps.
\item \textbf{Text Temporal LR}: Same temporal LR on sentence-level text embeddings (all-MiniLM-L6-v2).
\item \textbf{GRU Sequence Model}: Gated recurrent unit on hidden-state sequences across multiple layers.
\item \textbf{Mahalanobis Distance}: TrajGuard-style Mahalanobis distance on safe-trace covariance.
\item \textbf{BOCPD / CUSUM}: Bayesian Online Change-Point Detection and CUSUM on probe score sequences.
\item \textbf{16-Label Taxonomy}: Probes trained with fine-grained safety labels (4 types $\times$ pre/post).
\end{enumerate}

\subsection{Layer Selection Sensitivity Analysis}
\label{sec:layer_selection}

We evaluate all layers under five selection metrics:

\begin{enumerate}
\item \texttt{bal\_acc $\times$ crossing\_rate}: Composite metric penalizing low crossing rates.
\item \texttt{pure bal\_acc}: Threshold-free balanced accuracy.
\item \texttt{precision}: Threshold-dependent positive predictive value.
\item \texttt{pre\_cp\_fpr}: Threshold-dependent pre-commitment-point false positive rate.
\item \texttt{composite (additive)}: Weighted sum of precision, inverse FPR, and coverage.
\end{enumerate}

For each metric, we select the optimal layer per model, then evaluate downstream HS-vs-text comparisons. This design exposes how metric choice determines which layer is ``best'' and whether the HS advantage is visible.

\subsection{Statistical Framework}

All comparisons use 10,000-replicate bootstrap with Holm--Bonferroni correction across 12 tests (4 models $\times$ 3 metrics: balanced accuracy, precision, FPR). We report 95\% confidence intervals and adjusted $p$-values.


\section{Experiments and Results}

\subsection{Layer Selection Bias}
\label{sec:layer_bias}

Table~\ref{tab:layer_selection} presents selected layers across five metrics and four models. The biased metric \texttt{bal\_acc $\times$ crossing\_rate} selects extremely shallow layers for 32B models: L2 (3\% depth) for QwQ-32B and L10 (16\%) for R1-32B. In contrast, threshold-dependent metrics (precision, pre\_cp\_fpr, composite) converge on deep layers: L60 (94\%) and L63 (98\%) respectively.

\begin{table}[t]
\centering
\small
\begin{tabular}{lcccc}
\toprule
\textbf{Metric} & \textbf{R1-8B} & \textbf{OT-7B} & \textbf{QwQ} & \textbf{R1-32B} \\
\midrule
bal\_acc $\times$ cr & L0 (0\%) & L0 (0\%) & L2 (3\%) & L10 (16\%) \\
pure bal\_acc & L12 (39\%) & L9 (32\%) & L30 (47\%) & L47 (73\%) \\
precision & L16 (52\%) & L16 (57\%) & L60 (94\%) & L63 (98\%) \\
pre\_cp\_fpr & L16 (52\%) & L16 (57\%) & L60 (94\%) & L63 (98\%) \\
composite & L16 (52\%) & L12 (43\%) & L60 (94\%) & L63 (98\%) \\
\bottomrule
\end{tabular}
\caption{Layer selected by each metric (depth \%) across four models. The biased composite \texttt{bal\_acc $\times$ crossing\_rate} selects shallow layers for 32B models (3--16\% depth), while three threshold-dependent metrics converge on deep layers (94--98\%).}
\label{tab:layer_selection}
\end{table}

Under leave-one-out cross-validation, the dominant layer is selected in 65--68\% of folds regardless of metric, confirming selection stability. The \texttt{crossing\_rate} factor biases toward shallow layers because compressed early representations produce more frequent threshold crossings (noise-driven), inflating this term for layers where the probe signal is weakest.

\paragraph{Impact on conclusions.}
At biased shallow layers, the HS-vs-text precision gap is non-significant: QwQ-32B L2 yields $+$5.8pp ($p=0.35$), R1-32B L10 yields $+$9.4pp ($p=0.064$). At corrected deep layers, the same comparison produces highly significant results (Table~\ref{tab:hb_summary}).

\subsection{HS vs.\ Text at Corrected Layers}
\label{sec:hs_vs_text}

Table~\ref{tab:hb_summary} presents the main HS-vs-text comparison at corrected layers with Holm--Bonferroni correction.

\begin{table}[t]
\centering
\small
\begin{tabular}{lccccc}
\toprule
\textbf{Model} & \textbf{Layer} & \textbf{W} & \textbf{$\Delta$Prec} & \textbf{$\Delta$BA} & \textbf{$\Delta$FPR} \\
\midrule
R1-8B & L14 & 15 & +22.4\rlap{**} & +8.2\rlap{**} & $-$7.4\rlap{*} \\
OT-7B & L16 & 3 & +21.9\rlap{**} & +10.7\rlap{**} & $-$17.1\rlap{**} \\
QwQ-32B & L60 & 3 & +23.6\rlap{**} & +15.4\rlap{**} & $-$9.8\rlap{**} \\
R1-32B & L63 & 1 & +15.4\rlap{**} & +18.7\rlap{**} & $-$14.6\rlap{**} \\
\bottomrule
\end{tabular}
\caption{HS vs.\ text precision ($\Delta$Prec), balanced accuracy ($\Delta$BA), and false positive rate ($\Delta$FPR) deltas in percentage points at corrected layers. All 12/12 tests pass Holm--Bonferroni at $\alpha{=}0.05$. ** $p{<}0.01$, * $p{<}0.05$ (HB-adjusted). Note: window sizes differ per model.}
\label{tab:hb_summary}
\end{table}

All 12 tests pass the step-down procedure. Precision advantages are remarkably consistent ($+$15.4 to $+$23.6pp). However, three caveats qualify this result: (1) window sizes differ per model (W=1 for R1-32B, W=3 for OT-7B/QwQ, W=15 for R1-8B); (2) R1-32B at W=1 achieves only 25\% coverage (5/20 traces detected); (3) QwQ-32B HS lead time is negative ($-$1.2 to $+$1.9 steps) while text detects 3.7--12.3 steps earlier.

\paragraph{BA-layer robustness check.}
To address circularity (selecting by precision, evaluating by precision), we repeat the comparison at balanced-accuracy-optimal layers (threshold-free criterion). All 16/16 tests (4 models $\times$ 4 metrics) pass Holm--Bonferroni, confirming HS advantage is a modality property, not a layer selection artifact. Notably, R1-32B at BA-optimal L47 shows $\Delta$Prec $+$13.7pp and $\Delta$AUROC $+$21.1pp.

\subsection{Temporal Window Effects}
\label{sec:temporal}

Table~\ref{tab:temporal} shows the effect of temporal context on HS probe accuracy at corrected layers.

\begin{table}[t]
\centering
\small
\begin{tabular}{llccccc}
\toprule
\textbf{Model} & \textbf{Layer} & \textbf{W=1} & \textbf{W=3} & \textbf{W=5} & \textbf{W=10} & \textbf{W=15} \\
\midrule
R1-32B & L63 & 0.826* & 0.754 & 0.741 & 0.737 & 0.721 \\
QwQ-32B & L60 & 0.740 & 0.740 & 0.742 & 0.735 & 0.731 \\
R1-8B & L14 & 0.778 & 0.790 & 0.781 & 0.793 & 0.802 \\
OT-7B & L16 & 0.663 & 0.688* & 0.678 & 0.668 & 0.659 \\
\bottomrule
\end{tabular}
\caption{Balanced accuracy of HS temporal probes by window size. R1-8B improves with larger $W$; R1-32B \emph{degrades} ($-$10.5pp by W=15) due to dimensionality constraints (51,200 $\times$ $W$ features for $n{=}20$). * marks best $W$ per model.}
\label{tab:temporal}
\end{table}

Text probes improve universally with larger windows across all four models. HS temporal improvement is model-dependent: R1-8B benefits monotonically (0.778$\to$0.802), OT-7B peaks at W=3 (0.688) then declines, QwQ-32B is flat, and R1-32B degrades sharply due to the curse of dimensionality ($51{,}200 \times W$ features with only $n{=}20$ samples).

\paragraph{Precision--coverage tradeoff.}
Larger windows universally increase precision and decrease coverage. At the best window per model, coverage ranges from 48.8\% (R1-8B, W=25, precision 0.822) to 70.0\% (R1-32B, W=10, precision 0.609). No operating point achieves $>$0.70 precision and $>$0.60 coverage simultaneously.

\subsection{Alternative Methods: Comprehensive Negative Results}

\paragraph{Mahalanobis distance.}
TrajGuard-style Mahalanobis distance fails on all four models. Table~\ref{tab:mahal} shows that safe vs.\ post-commitment separation is negligible (max $+$0.20 on QwQ-32B), with false alarms exceeding true positives at every threshold.

\begin{table}[t]
\centering
\small
\begin{tabular}{lcccc}
\toprule
\textbf{Model} & \textbf{Layer} & \textbf{Safe MD} & \textbf{Post-CP MD} & \textbf{Sep.} \\
\midrule
R1-8B & L14 & 6.66 & 6.73 & +0.07 \\
OT-7B & L16 & 7.68 & 7.71 & +0.03 \\
R1-32B & L63 & 6.73 & 6.88 & +0.15 \\
QwQ-32B & L60 & 3.22 & 3.42 & +0.20 \\
\bottomrule
\end{tabular}
\caption{Mahalanobis distance from safe-trace covariance. Separation is $<$0.21 for all models; false alarm rate exceeds true positive rate at every threshold.}
\label{tab:mahal}
\end{table}

\paragraph{BOCPD and CUSUM.}
BOCPD detects commitment on only 1/4 models (R1-8B, 93\% detection, lead $+$3.45 steps). On the remaining three models, detection rate is 0\%. CUSUM detects on 3/4 models but achieves worse lead time than simple first-crossing-point (FCP) thresholds on all models (Table~\ref{tab:cpd}).

\begin{table}[t]
\centering
\small
\begin{tabular}{lccccc}
\toprule
\textbf{Model} & \textbf{BOCPD \%} & \textbf{BOCPD Lead} & \textbf{CUSUM \%} & \textbf{FCP \%} & \textbf{FCP Lead} \\
\midrule
R1-8B & 93 & +3.45 & 65 & 65 & +4.89 \\
OT-7B & 0 & --- & 75 & 89 & +11.83 \\
R1-32B & 0 & --- & 85 & 95 & +5.68 \\
QwQ-32B & 0 & --- & 74 & 79 & +9.40 \\
\bottomrule
\end{tabular}
\caption{Change-point detection methods. BOCPD fails on 3/4 models. FCP (simple threshold) achieves comparable or better detection and lead time than both BOCPD and CUSUM.}
\label{tab:cpd}
\end{table}

\paragraph{GRU sequence model.}
GRU does not outperform temporal LR on balanced accuracy for 3/4 models ($-$6.2pp on R1-32B, $-$3.7pp on QwQ-32B, $+$0.3pp on R1-8B). On OT-7B, GRU achieves $+$5.8pp over LR but at the cost of 2.7$\times$ higher FPR (0.274 vs.\ 0.100). GRU trades accuracy for coverage and lead time via higher false positive rates (Table~\ref{tab:gru}).

\begin{table}[t]
\centering
\small
\begin{tabular}{lcccc}
\toprule
\textbf{Model} & \textbf{LR BA} & \textbf{GRU BA} & \textbf{$\Delta$} & \textbf{GRU FPR} \\
\midrule
R1-32B & 0.737 & 0.675 & $-$6.2 & 0.135 \\
QwQ-32B & 0.761 & 0.724 & $-$3.7 & 0.289 \\
OT-7B & 0.688 & 0.746 & +5.8 & 0.274 \\
R1-8B & 0.806 & 0.809 & +0.3 & 0.165 \\
\bottomrule
\end{tabular}
\caption{GRU vs.\ temporal LR. GRU improves over LR on only 1/4 models (OT-7B), confirming linear probes are at or near the information ceiling for this task.}
\label{tab:gru}
\end{table}

\paragraph{16-label taxonomy.}
Fine-grained supervision (16 labels) provides marginal improvement after binary collapse: $+$0.6 to $+$1.1pp across models (Table~\ref{tab:taxonomy}). Raw 16-class accuracy is low (26--41\%), indicating that transition subtypes are not cleanly separable in hidden-state space.

\begin{table}[t]
\centering
\small
\begin{tabular}{lcccc}
\toprule
\textbf{Model} & \textbf{Binary} & \textbf{16-Label} & \textbf{16$\to$Bin} & \textbf{$\Delta$} \\
\midrule
R1-8B & 0.817 & 0.41 & 0.825 & +0.8 \\
OT-7B & 0.688 & 0.35 & 0.699 & +1.1 \\
QwQ-32B & 0.740 & 0.26 & 0.749 & +0.9 \\
R1-32B & 0.826 & 0.28 & 0.832 & +0.6 \\
\bottomrule
\end{tabular}
\caption{16-label vs.\ binary supervision. After collapsing 16 labels to binary, improvement is marginal ($\leq$1.1pp), suggesting linear probes are saturated at the binary detection boundary.}
\label{tab:taxonomy}
\end{table}

\subsection{Cross-Model Transfer}

Probes trained on one model and tested on another fail comprehensively (Table~\ref{tab:transfer}). The best cross-model accuracy is 67.3\% (within the Qwen-32B family), far below within-model performance and barely above the 50\% chance baseline. Cross-family transfer (Llama$\leftrightarrow$Qwen) is at chance.

\begin{table}[t]
\centering
\small
\begin{tabular}{lcccc}
\toprule
\textbf{Train $\to$} & \textbf{R1-8B} & \textbf{OT-7B} & \textbf{R1-32B} & \textbf{QwQ} \\
\midrule
R1-8B & \textit{0.817} & --- & 0.522 & 0.595 \\
OT-7B & --- & \textit{0.688} & --- & --- \\
R1-32B & 0.583 & --- & \textit{0.826} & 0.673 \\
QwQ-32B & --- & --- & 0.673 & \textit{0.740} \\
\bottomrule
\end{tabular}
\caption{Cross-model transfer accuracy. Self-accuracy (diagonal, italic) far exceeds any transfer pair. The best transfer (67.3\%, Qwen-32B internal) is 7--15pp below self-accuracy. Cross-family transfer is at chance.}
\label{tab:transfer}
\end{table}

\subsection{Lead Time Analysis}

Table~\ref{tab:lead} compares lead time---steps between first detection and commitment point---across methods and models. Results are mixed: R1-8B HS detects \emph{at} the commitment point (lead $\approx$0), QwQ-32B HS has negative lead ($-$1.2 to $+$1.9 steps), and only R1-32B shows substantial positive HS lead ($+$11.7 to $+$13.3 steps), but at only 25\% coverage.

\begin{table}[t]
\centering
\small
\begin{tabular}{lcccc}
\toprule
\textbf{Model} & \textbf{HS Lead} & \textbf{Text Lead} & \textbf{HS BA} & \textbf{Pos BA} \\
\midrule
R1-8B (W=15) & $-$0.1 & +10.3$\dagger$ & 0.802 & --- \\
OT-7B (W=3) & +21.6$\ddagger$ & +45.5$\ddagger$ & 0.688 & --- \\
QwQ-32B (W=3) & $-$1.2 to +1.9 & +3.7 to +12.3 & 0.740 & 0.746 \\
R1-32B (W=1) & +11.7 to +13.3 & --- & 0.826 & 0.686 \\
\bottomrule
\end{tabular}
\caption{HS vs.\ text lead time. $\dagger$: inflated by high pre-CP FPR. $\ddagger$: inflated by long traces ($\sim$80 sentences). QwQ-32B position baseline (Pos BA) exceeds HS BA, questioning the value of white-box access for that model.}
\label{tab:lead}
\end{table}

\paragraph{Oracle decomposition.}
Oracle labels show gradual safety erosion pre-commitment ($+$27 to $+$31 step lead), while HS probes detect only at or near the commitment point. The 26--28 step gap between oracle and probe lead times indicates that hidden states encode \emph{current} safety state rather than the \emph{trajectory} toward commitment. On R1-8B, oracle label autocorrelation is 0.82 pre-CP (gradual transition) while HS entropy is 0.41 (step-function).

\subsection{False Positive Analysis}
\label{sec:fpr}

Table~\ref{tab:fpr_safe} reports false positive rates on safe traces---traces where the model never commits to harmful output. No method achieves simultaneously low FPR, high precision, and high coverage.

\begin{table}[t]
\centering
\small
\begin{tabular}{lcccc}
\toprule
\textbf{Model} & \textbf{HS Static} & \textbf{HS Temp} & \textbf{Text Temp} & \textbf{GRU} \\
\midrule
R1-8B & 0.390 & 0.392 & 0.385 & 0.521 \\
OT-7B & 0.111 & 0.175 & 0.460 & --- \\
QwQ-32B & 0.500 & 0.450 & 0.425 & --- \\
R1-32B & 0.239 & 0.239 & 0.310 & --- \\
\bottomrule
\end{tabular}
\caption{Safe-trace false positive rates. Median FPR $\approx$0.39 across methods. QwQ-32B HS FPR (0.45--0.50) \emph{exceeds} text FPR (0.325--0.425).}
\label{tab:fpr_safe}
\end{table}

\paragraph{Deployment viability.}
We define three practical criteria: safe FPR $<$20\%, precision $>$70\%, and coverage $>$60\%. Only 1/4 models meets each criterion individually (OT-7B for FPR, R1-8B for precision at W$\geq$15, R1-8B/R1-32B for coverage at small W). Zero models meet all three simultaneously.

\subsection{Bootstrap Confidence Intervals}
\label{sec:bootstrap}

Table~\ref{tab:hb_full} presents the full Holm--Bonferroni step-down procedure.

\begin{table}[t]
\centering
\small
\begin{tabular}{clccc}
\toprule
\textbf{Rank} & \textbf{Model/Metric} & \textbf{$p$} & \textbf{HB Thr.} & \textbf{Pass} \\
\midrule
1 & R1-8B / precision & $<$.001 & .0042 & \checkmark \\
2 & OT-7B / bal\_acc & $<$.001 & .0045 & \checkmark \\
3 & OT-7B / precision & $<$.001 & .0050 & \checkmark \\
4 & OT-7B / FPR & $<$.001 & .0056 & \checkmark \\
5 & QwQ / bal\_acc & $<$.001 & .0063 & \checkmark \\
6 & QwQ / precision & $<$.001 & .0071 & \checkmark \\
7 & R1-32B / bal\_acc & $<$.001 & .0083 & \checkmark \\
8 & R1-32B / precision & $<$.001 & .0100 & \checkmark \\
9 & QwQ / FPR & .004 & .0125 & \checkmark \\
10 & R1-8B / bal\_acc & .005 & .0167 & \checkmark \\
11 & R1-32B / FPR & .007 & .0250 & \checkmark \\
12 & R1-8B / FPR & .044 & .0500 & \checkmark \\
\bottomrule
\end{tabular}
\caption{Holm--Bonferroni step-down results. 12/12 pass ($\alpha{=}0.05$). Tests ranked by $p$-value. CI widths for 32B models are wide ($\sim$0.12--0.15 for BA) due to small $n$.}
\label{tab:hb_full}
\end{table}

\subsection{Ensemble Cross-Validation (7--8B)}

Five-fold cross-validated ensembles of metric-weighted layer selections on 7--8B models confirm HS advantage over text (Table~\ref{tab:ensemble}). Ensemble precision is lower than the best single metric on R1-8B (0.565 vs.\ 0.604) but slightly higher on OT-7B (0.506 vs.\ 0.473).

\begin{table}[t]
\centering
\small
\begin{tabular}{lcccc}
\toprule
\textbf{Model} & \textbf{Ens.\ Prec} & \textbf{Best Single} & \textbf{Text Prec} & \textbf{Ens vs.\ Text $p$} \\
\midrule
R1-8B & 0.565 & 0.604 & 0.372 & $<$0.002 \\
OT-7B & 0.506 & 0.473 & 0.236 & $<$0.001 \\
\bottomrule
\end{tabular}
\caption{Cross-validated ensemble results on 7--8B models. HS advantage over text is validated under CV ($p \leq 0.002$).}
\label{tab:ensemble}
\end{table}


\section{Discussion}
\label{sec:discussion}

\paragraph{Layer selection as researcher degree of freedom.}
Our central finding is that layer selection metric choice is a potent confound in hidden-state safety evaluation. The composite \texttt{bal\_acc $\times$ crossing\_rate} selects shallow layers for 32B models (3--16\% depth), where probe signal is weakest, entirely masking the HS advantage. Three of four non-biased metrics converge on deep layers (94--98\%), while threshold-free balanced accuracy selects middle layers (47--73\%). This divergence constitutes a second axis of researcher degrees of freedom beyond the crossing-rate bias. We recommend that future work report HS evaluation results under \emph{at least} three selection criteria spanning threshold-dependent and threshold-free metrics.

\paragraph{Cost-benefit analysis.}
HS probes yield $+$15--24pp precision over text probes. This advantage requires: (a) local model hosting (no API access), (b) inference pipeline modification for intermediate activation extraction, (c) per-model layer selection and probe training. Text probes require only access to the generated CoT---available in any deployment. At text precision of 0.33--0.49 (1 true alarm per 2--3 triggers) vs.\ HS precision of 0.47--0.78 (1 per 1.3--2.1 triggers), the marginal improvement may not justify white-box infrastructure costs for many deployment settings.

\paragraph{Fundamental limitation: reactive, not predictive.}
Oracle decomposition reveals a 26--28 step gap between what text labels encode (gradual erosion) and what HS probes detect (step-function at commitment). Hidden states encode the model's \emph{current} safety state, not the trajectory toward commitment. For deployment, trajectory prediction provides earlier warning, suggesting text-based monitoring of reasoning patterns---while noisier---may be the more promising research direction for early detection.

\paragraph{Cross-model transfer failure.}
Complete transfer failure (best 67.3\%) implies that safety commitment encoding is model-specific. Every new model requires its own probe training with labeled commitment points, a supervision cost that limits practical scalability. Combined with the engineering cost of white-box access, the deployment case for HS monitoring is substantially weakened.


\section{Conclusion}

We present the first systematic evaluation of safety commitment detection in reasoning models, comparing seven methods across four models. Our primary contribution is the discovery that layer selection metric choice is a critical confounder: the composite metric \texttt{bal\_acc $\times$ crossing\_rate} systematically biases toward shallow layers at 32B scale, reversing conclusions about hidden-state monitoring efficacy. After correction, HS probes significantly outperform text probes (12/12 Holm--Bonferroni, precision $+$15--24pp), but this statistical advantage must be weighed against practical limitations: no method achieves deployment-ready performance (0/4 models meet all practical criteria), safe-trace FPR is 39\% at the median, cross-model transfer fails, and hidden states are reactive rather than predictive. We release our multi-metric audit framework and urge the community to adopt multi-criterion layer selection reporting as standard practice.


\section*{Limitations}

Our evaluation has three principal limitations. First, test set sizes for 32B models ($n{=}19$--20) yield wide confidence intervals; the R1-32B residualized HS-text gap ($+$1.9pp) is below the minimum detectable effect (11.2pp at 80\% power), so the non-significant result reflects a power limitation rather than evidence of absence. Second, the text baseline uses a single encoder (all-MiniLM-L6-v2); stronger encoders or LLM-based classifiers would likely narrow the HS-text gap. Third, all probes are trained on commitment-point annotations that themselves require human judgment, introducing a potential source of label noise.


\bibliographystyle{acl_natbib}
\begin{thebibliography}{10}

\bibitem[\protect\citename{DeepSeek-AI}2025]{deepseekai2025deepseekr1}
DeepSeek-AI. 2025.
\newblock DeepSeek-R1: Incentivizing reasoning capability in LLMs via reinforcement learning.
\newblock {\em arXiv preprint arXiv:2501.12948}.

\bibitem[\protect\citename{Qwen Team}2025]{qwq32b}
Qwen Team. 2025.
\newblock QwQ: Reflect deeply on the boundaries of the unknown.
\newblock Blog post.

\bibitem[\protect\citename{Wen et al.}2025]{wen2025thinkingmakesllmssusceptible}
Sijia Wen, Haoran Li, and Yixin Cao. 2025.
\newblock Thinking makes LLMs susceptible to jailbreaking.
\newblock {\em arXiv preprint}.

\bibitem[\protect\citename{Zhou et al.}2025]{zhou2025hiddenriskslargereasoning}
Yihe Zhou, et al. 2025.
\newblock Hidden risks of large reasoning models.
\newblock {\em arXiv preprint}.

\bibitem[\protect\citename{Wang et al.}2024]{wang2024mitigating}
Zhixin Wang, et al. 2024.
\newblock Mitigating jailbreak attacks via trajectory-aware guard.
\newblock {\em arXiv preprint}.

\bibitem[\protect\citename{Ji et al.}2024]{ji2024llmguard}
Ziwei Ji, et al. 2024.
\newblock LLM guard: Defending LLMs against jailbreaks via internal state monitoring.
\newblock {\em arXiv preprint}.

\bibitem[\protect\citename{Li et al.}2024]{li2024inferencetime}
Kenneth Li, et al. 2024.
\newblock Inference-time intervention: Eliciting truthful answers from a language model.
\newblock {\em NeurIPS 2024}.

\bibitem[\protect\citename{Burns et al.}2023]{burns2023discovering}
Collin Burns, et al. 2023.
\newblock Discovering latent knowledge in language models without supervision.
\newblock {\em ICLR 2023}.

\bibitem[\protect\citename{Zou et al.}2023]{zou2023representation}
Andy Zou, et al. 2023.
\newblock Representation engineering: A top-down approach to AI transparency.
\newblock {\em arXiv preprint}.

\bibitem[\protect\citename{Belrose et al.}2023]{belrose2023eliciting}
Nora Belrose, et al. 2023.
\newblock Eliciting latent predictions from transformers with the tuned lens.
\newblock {\em arXiv preprint}.

\bibitem[\protect\citename{Zheng et al.}2025]{zheng2025safethink}
Zheng Zheng, et al. 2025.
\newblock SafeThink: Safe thinking for LLM reasoning.
\newblock {\em arXiv preprint}.

\end{thebibliography}

\end{document}
